\begin{frame}{微观模型：分子 Hamiltonian 与腔场偶极耦合}
  \setlength{\abovedisplayskip}{0pt}
  \setlength{\belowdisplayskip}{0pt}
  \[
    H=H_0+V,
    \quad
    H_0=H_{\mathrm{ph}}+H_{\mathrm{mol}},
    \quad
    H_{\mathrm{ph}}=\hbar\omega_{\mathrm{ph}}a^\dagger a,
  \]

  \[
    V=-\physmark{Plum}{p4lambda}{\hbar\lambda}
    (a+a^\dagger)
    \physmark{Bittersweet}{p4mu}{\hat\mu},
    \quad
    \hat\mu=\sum_{i=1}^{N}\hat\mu_i(\bm q_i,\bm Q_i),
    \quad
    \hbar\lambda=
    \sqrt{\tfrac{\hbar\omega_{\mathrm{ph}}}{2\epsilon_0V}}.
  \]
  \begin{tikzpicture}[overlay,remember picture,>=stealth,
      nodes={align=center,inner ysep=1pt,font=\scriptsize},<-]
    \path (p4lambda.south) ++ (-1.5,-0.4em)
      node[anchor=north,color=Plum!85] (p4lambdatext)
      {\textsf{\footnotesize 真空腔场耦合强度}};
    \draw[color=Plum] (p4lambda.south) |- (p4lambdatext.north);
    \path (p4mu.south) ++ (1.5,-1em)
      node[anchor=north,color=Bittersweet!85] (p4mutext)
      {\textsf{\footnotesize 分子系综总偶极矩}};
    \draw[color=Bittersweet] (p4mu.south) |- (p4mutext.north);
  \end{tikzpicture}

  \[
    H_{\mathrm{mol}}
    =\sum_{i=1}^{N}
      \physmark{sintefgreen}{p3Hi}{H_i}
      \bigl(\physmark{OliveGreen}{p3qi}{\bm q_i},
      \physmark{sintefdarkgreen}{p3Qi}{\bm Q_i}\bigr)
    =\sum_y
      \physmark{sintefdarkgreen}{p3ey}{\hbar\Omega_y|y\rangle\langle y|}.
  \]
  \begin{tikzpicture}[overlay,remember picture,>=stealth,
      nodes={align=center,inner ysep=1pt,font=\scriptsize},<-]
    \path (p3qi.south) ++ (-0.55,-1.5em)
      node[anchor=north,color=OliveGreen!85] (p3qitext)
      {\textsf{\footnotesize 电子}};
    \draw[color=OliveGreen] (p3qi.south) |- (p3qitext.north);
    \path (p3Qi.south) ++ (0.55,-1.5em)
      node[anchor=north,color=sintefdarkgreen!85] (p3Qitext)
      {\textsf{\footnotesize 核}};
    \draw[color=sintefdarkgreen] (p3Qi.south) |- (p3Qitext.north);
    \path (p3ey.south) ++ (0.85,-1.5em)
      node[anchor=north,color=sintefdarkgreen!85] (p3eytext)
      {\textsf{\footnotesize 本征能量与多体本征态}};
    \draw[color=sintefdarkgreen] (p3ey.south) |- (p3eytext.north);
  \end{tikzpicture}
  \vspace{1.2em}

  {\centering\scriptsize
    电子结构、核振动、非谐性与多能级跃迁仍保留在
    $H_i(\bm q_i,\bm Q_i)$ 和 $\hat\mu_i$ 中。\par}
  \papersource{Yuen-Zhou \& Koner, JCP 160, 154107 (2024), Eqs.~(1)--(5).}
\end{frame}

\begin{frame}[c]{相同二能级分子的集体亮态与暗态}
  \small
  $N$ 个相同二能级分子、均匀单分子耦合 $g$、单激发子空间 定义亮态与暗态：

  \vspace{0.2em}
  {\centering\normalsize
    \(\displaystyle
      |B\rangle=\frac{1}{\sqrt N}\sum\nolimits_{j=1}^{N}|e_j\rangle,
    \qquad
      \langle D_\mu|B\rangle=0
    \)
    \par}

  \vspace{0.2em}
  {
    集体亮态是单激发在 $N$ 个分子之间的对称相干叠加；在相同二能级分子与旋波近似下

    ( Tavis--Cummings):\par}

  \vspace{1.5em}

  {\centering\small
    \(\displaystyle
      \begin{aligned}
        V_{\mathrm{RWA}}
          &=\hbar g\sum\nolimits_{j=1}^{N}
            \left(a^\dagger\sigma_j^-+a\sigma_j^+\right),
        \qquad
        B^\dagger
          =\frac{1}{\sqrt N}\sum\nolimits_{j=1}^{N}\sigma_j^+,
        \\ \\
        V_{\mathrm{RWA}}
          &=\hbar\physmark{Plum}{p5Gcoupling}{G}
            \left(a^\dagger B+aB^\dagger\right),
        \qquad
        \physmark{Plum}{p5G}{G}
          =\physmark{NavyBlue}{p5g}{g}
            \physmark{OliveGreen}{p5sqrtN}{\sqrt N}.
      \end{aligned}
    \)
    \par}
  \begin{tikzpicture}[overlay,remember picture,>=stealth,
      nodes={align=center,inner ysep=1pt},<-]
    \path (p5G.north) ++ (-1.25,-2.6em)
      node[anchor=east,color=Plum!85] (p5Gtext)
      {\textsf{\footnotesize 集体耦合}};
    \draw[color=Plum] (p5G.south) |- (p5Gtext.east);
    \path (p5g.north) ++ (0,-3.6em)
      node[anchor=south,color=NavyBlue!85] (p5gtext)
      {\textsf{\footnotesize 单分子耦合}};
    \draw[color=NavyBlue] (p5g.south) -- (p5gtext.north);
    \path (p5sqrtN.north) ++ (1.05,-2.6em)
      node[anchor=west,color=OliveGreen!85] (p5sqrtNtext)
      {\textsf{\footnotesize 集体相干增强}};
    \draw[color=OliveGreen]
      (p5sqrtN.south) |- (p5sqrtNtext.west);
  \end{tikzpicture}

  \vspace{2.0em}
\end{frame}

\begin{frame}[c]{腔光子与集体亮态杂化形成分子极化激元}
  \small
  集体亮态 $|B\rangle$与腔光子态构成基矢
  $|1_{\mathrm{cav}},\mathrm{GS}\rangle$、$|0_{\mathrm{cav}},B\rangle$。

  \vspace{0.8em}
  {\large
  \begin{equation*}
    H_{\mathrm{1exc}}=\hbar
    \begin{pmatrix}
      \physmark{Plum}{p6wc}{\omega_c} &
      \physmark{Bittersweet}{p6G}{G}\\[0.2em]
      G &
      \physmark{NavyBlue}{p6w0}{\omega_0}
    \end{pmatrix}.
  \end{equation*}}

  \begin{tikzpicture}[overlay,remember picture,>=stealth,
      nodes={align=center,inner ysep=1pt},<-]

    \path (p6wc.north) ++(-1.3,1.1em)
      node[anchor=south,color=Plum!85] (p6wctext)
      {\textsf{\footnotesize 单模腔频率}};
    \draw[color=Plum] (p6wc.north) |- (p6wctext.south);

    \path (p6w0.south) ++(1.4,-1.5em)
      node[anchor=north,color=NavyBlue!85] (p6w0text)
      {\textsf{\footnotesize 集体亮态频率}};
    \draw[color=NavyBlue] (p6w0.south) |- (p6w0text.north);

    \path (p6G.east) ++(1.4,1.1em)
      node[anchor=west,color=Bittersweet!85] (p6Gtext)
      {\textsf{\footnotesize 集体耦合}\\[-0.1em]
       \textsf{\footnotesize $G=g\sqrt{N}$}};
    \draw[color=Bittersweet] (p6G.north) |- (p6Gtext.west);

  \end{tikzpicture}

  \vspace{1.4em}

  {\large
  \begin{equation*}
    \omega_\pm=\frac{\omega_c+\omega_0}{2}
    \pm\sqrt{G^2+\frac{(\omega_c-\omega_0)^2}{4}}.
  \end{equation*}}

  \vspace{0.4em}

  {\footnotesize
    分子极化激元：腔光子模式与分子集体激发强耦合杂化的光--物质混合本征模（UP/LP）。\par}

\end{frame}

\begin{frame}[c]{分子极化与线性极化率}
  \small
  在线性响应近似下，只保留外场诱导极化的一阶项：

  {\centering\large
    \[
      \physmark{Plum}{p7dP}{\delta P^{(1)}(t)}
      =
      \int_{-\infty}^{t}
      \physmark{NavyBlue}{p7chi1}{\chi^{(1)}(t-t')}
      \physmark{Bittersweet}{p7Eweak}{E(t')}\,dt' .
    \]
    \par
  }

  \begin{tikzpicture}[
      overlay,
      remember picture,
      >=stealth,
      nodes={align=center,inner ysep=1pt},
      <-
    ]

    \path (p7dP.south) ++(-1.15,-0.75em)
      node[
        anchor=north,
        color=Plum!85
      ] (p7dPtext)
      {\textsf{\footnotesize 一阶诱导极化}};

    \draw[color=Plum]
      (p7dP.south) |- (p7dPtext.north);

    \path (p7chi1.south) ++(0,-0.75em)
      node[
        anchor=north,
        color=NavyBlue!85
      ] (p7chi1text)
      {\textsf{\footnotesize 线性极化响应}};

    \draw[color=NavyBlue]
      (p7chi1.south) |- (p7chi1text.north);

    \path (p7Eweak.south) ++(1.05,-0.75em)
      node[
        anchor=north,
        color=Bittersweet!85
      ] (p7Eweaktext)
      {\textsf{\footnotesize 弱外电场}};

    \draw[color=Bittersweet]
      (p7Eweak.south) |- (p7Eweaktext.north);

  \end{tikzpicture}

  \vspace{1.4em}

  \small
  $\chi^{(1)}(t-t')$ 描述分子受到时刻 $t'$ 的微弱电场扰动后，
  在时刻 $t$ 保留下来的偶极响应。

  {\centering\large
  \[
    \delta P^{(1)}(\omega)
    =
    \chi^{(1)}(\omega)E(\omega).
  \]
  \par}

  \vspace{0.5em}
  \small
  $\chi^{(1)}(\omega)$ 表征分子对频率为 $\omega$ 的弱外场产生的
  线性偶极响应强度与相位。

  \vspace{0.3em}
  {\color{gray}
    $\chi(\omega)$ 可由微观动力学或线性响应理论求得又或者线性吸收、色散或复折射率实验。
  }

\end{frame}
